A set \(A\subseteq\mathbb{N}\) is \emph{dissociated} if every relation
\[
\sum_{i=1}^t \varepsilon_i a_i = 0,\qquad a_i\in A\ \text{distinct},\ \varepsilon_i\in\{-1,0,1\},
\]
forces all \(\varepsilon_i=0\).
(For positive integers, this is equivalent to uniqueness of subset sums.)

Say that \(A\) is \emph{proportionately dissociated} if there exists \(\delta>0\) such that every finite \(B\subseteq A\) contains a dissociated \(C\subseteq B\) with \(|C|\ge \delta|B|\).

Assume additionally that \(A\) is a union of finitely many dissociated sets.
Question: must \(A\) in fact be a union of \emph{at most} \(M(\delta)\) dissociated sets for some bound depending only on \(\delta\)?
